% =====================================================================
%  CONJURA CONJECTURE STATEMENT
%  Breaking fully quantum time-lock puzzles in the QROM.
%  Requires conjura-conjecture.cls in the same directory.
% =====================================================================
\documentclass{conjura-conjecture}

\runninghead{FULLY QUANTUM TIME-LOCK PUZZLES IN THE QROM}

\newcommand{\Gen}{\mathsf{Gen}}
\newcommand{\Sol}{\mathsf{Sol}}
\newcommand{\Ext}{\mathsf{E}}
\newcommand{\puz}{\mathsf{P}}
\newcommand{\solu}{\mathsf{s}}
\newcommand{\solhon}{\mathsf{s}'}
\newcommand{\solatt}{\mathsf{s}''}

\begin{document}

\cjkicker{CONJURA \textperiodcentered{} OPEN PROBLEM}
\cjtitle{Breaking Fully Quantum Time-Lock Puzzles in the Quantum Random
  Oracle Model}
\cjsubtitle{Both the puzzle generator and the honest solver are quantum,
  and so is the attacker}
\cjstatus{Statement: AI-written from Abtin Afshar, Kai-Min Chung,
  Yao-Ching Hsieh, Yao-Ting Lin and Mohammad Mahmoody, \emph{On the
  (Im)possibility of Time-Lock Puzzles in the Quantum Random Oracle
  Model}, Cryptology ePrint Archive 2023 (no venue or paper number is
  printed on the PDF itself), not yet checked by a human. Settled in that
  paper: a classical generator with a quantum solver (for arbitrary
  completeness error, and round-optimally under perfect completeness), and
  a quantum generator with a classical solver (under perfect
  completeness); open, and conjectured below, is the fully quantum case,
  where the paper presents no attack and shows that any attack asking only
  classical queries would prove the Aaronson--Ambainis simulation
  conjecture.}
\cjcategory{\emph{Idealized Models \& Non-Uniformity}.}

\vspace{0.4em}

\begin{informalconjecture}
A time-lock puzzle locks a secret into a puzzle so that unwrapping it takes
a long time even for someone with many computers running in parallel, while
making the puzzle is cheap. If the only source of hardness is a random
oracle, one can measure ``time'' by the number of rounds of oracle queries
an algorithm needs, and ``parallelism'' by how many queries it may ask
inside one round. Classically such puzzles cannot exist: whatever the
generator does with its queries, an attacker recovers the solution using a
number of query rounds proportional to the number of queries the generator
made, no matter how many rounds the honest solver needs. The source paper
extends this impossibility to a quantum honest solver querying in
superposition, and separately to a quantum generator with a classical
solver. The remaining case is the one where generator and solver are
\emph{both} quantum; there the paper gives no attack, and it shows why its
technique stalls, by exhibiting a scheme --- based on a long-standing
conjecture about approximating quantum query algorithms classically being
false --- that no attacker restricted to classical queries can break even
with polynomially many queries and unlimited parallelism. The conjecture is
that the impossibility nevertheless holds once the attacker is allowed
quantum queries: for every such puzzle there is a quantum attacker that
recovers the solution about as well as the honest solver does, using
polynomially many quantum queries arranged into a number of rounds
proportional to the generator's query count.
\end{informalconjecture}

\section{The setting}\label{sec:setting}

Time-lock puzzles~\cite{RSW96} let a generator cheaply produce a puzzle
$\puz$ hiding a solution $\solu$ so that recovering $\solu$ takes a long
\emph{sequential} time, even for an adversary with unbounded parallelism.
Whether such puzzles can be built from symmetric-key primitives alone is
naturally studied in the random oracle model, where ``time'' has a clean
information-theoretic proxy: the number of \emph{rounds} of oracle queries
an algorithm needs, with arbitrarily many queries allowed inside a single
round.

In that model, Mahmoody, Moran and Vadhan~\cite{MMV11} ruled out time-lock
puzzles outright: if the generator asks $n$ oracle queries and the honest
solver asks $m$, there is always an attacker that recovers $\solu$ using
only $O(n)$ rounds of queries and $\poly(n,m)$ queries in total, however
sequential the honest solver is. They also showed $n$ rounds is the best one
can hope for, via a ``pseudo-chain'' puzzle that genuinely needs $n$ rounds.

Since random oracles can also be given in superposition~\cite{BDF11}, the
barrier splits into three settings, according to which of the generator and
the solver are quantum. The present paper closes two of them. When the
generator is classical and the solver quantum, it exhibits an attacker
asking $O(n)$ rounds of \emph{classical} queries and $\poly(n,m)$ queries
total, so the solver's quantum power cannot be what makes the puzzle hard; a
variant of the attack even runs in quantum polynomial time, which rules out
ROM constructions that lean on additional computational assumptions. When
the generator is quantum and the solver classical, it gives an attack in
exactly $n$ rounds and $m \cdot n$ queries for perfectly complete schemes,
by combining certificate-complexity ideas with the polynomial
representation of a quantum party's state used in the key-agreement
impossibility of~\cite{ACC22}. It also shows the pseudo-chain still needs
$n$ rounds against a solver making parallel \emph{quantum} queries, so
$O(n)$ rounds remains optimal in the quantum world.

The remaining case --- generator and solver both quantum --- is left open,
and the paper explains precisely what blocks the technique it uses
elsewhere. Every one of its attacks asks only classical queries, and
classical-query attacks cannot suffice unconditionally: assuming the
paper's asymptotic Quantum Polynomial-Query Simulation Conjecture
(Conjecture 6.1, p.~29) --- a weakening of the folklore simulation
conjecture formally stated as Conjecture 4 of Aaronson and
Ambainis~\cite{AA14} --- is \emph{false}, the paper builds a fully quantum
time-lock puzzle that no classical-query adversary breaks, even with
polynomially many queries and unlimited parallelism. It does so by
amplifying the conditional one-way-communication quantum key agreement
of~\cite{ACC22} to negligible completeness and soundness error and reading
the resulting transcript as a puzzle. The simulation conjecture --- that
the acceptance probability of an $n$-quantum-query algorithm can be
approximated on most oracles using $\poly(n)$ classical queries --- is
known only in the case where that probability is always $0$ or
$1$~\cite{OSSS05} and has resisted proof for over a decade, so the barrier
is unlikely to be removed soon. What is left, and what this conjecture asks
for, is an attack that itself uses quantum queries.

The stakes reach beyond time-lock puzzles. Under the reading of a
key-agreement transcript as a puzzle, observed in~\cite{MMV11} and used
here, round-efficient attacks on puzzles with quantum parties translate
into round-efficient attacks on key agreement in the quantum
ROM~\cite{ACC22}, and the paper notes that its techniques plausibly bear on
the impossibility of verifiable delay functions from random
oracles~\cite{MSW20}. Conversely, a resolution in the negative would be a
genuine positive result: proofs of sequential work \emph{are} achievable
against quantum adversaries from random oracles~\cite{CFHL21}, so the
boundary between what sequential primitives survive quantum access and what
do not is exactly what is at issue.

Progress to date. The paper settles the two mixed settings. Theorem 3.1
gives, for a classical generator and quantum solver, an attacker asking
$4n^{2}m^{2}/(\varepsilon\delta^{2})$ classical queries in $n/\varepsilon$
rounds with success at least $1 - \rho - \varepsilon - \delta$, and a second
variant with $2n$ rounds; Theorem 3.12 makes the attack run in quantum
polynomial time. Theorem 4.10 gives a round-optimal attack in at most $n$
rounds and at most $mn$ queries with success probability $1$ for perfectly
complete schemes with a classical generator and quantum solver, and Theorem
4.6 does the same for a quantum generator and classical solver. Theorem 5.3
shows the pseudo-chain puzzle of Mahmoody, Moran and Vadhan is still secure
against solvers making $q$ rounds of $k$-parallel quantum queries, up to
$O(\sqrt{kq^{3}/2^{\kappa}})$, so $O(n)$ rounds cannot be improved. For the
fully quantum case there is no attack, only the conditional barrier of
Theorem 6.2 against classical-query attackers.

\section{Notation and parameters}\label{sec:notation}

\begin{itemize}
  \item $\kappa \in \mathbb{N}$ is the security parameter. All asymptotics
    are in $\kappa$.
  \item $\mathcal{X}$ and $\mathcal{Y}$ are finite nonempty sets with
    $\mathcal{Y}$ carrying a group operation $\oplus$ (for concreteness one
    may take $\mathcal{X} = \mathcal{Y} = \{0,1\}^{\kappa}$ with bitwise
    XOR). The random oracle is $H \colon \mathcal{X} \to \mathcal{Y}$,
    drawn uniformly from the set of all such functions.
  \item Quantum access to $H$ is via the unitary
    $O_H \colon |x\rangle|y\rangle \mapsto |x\rangle|y \oplus H(x)\rangle$
    on $\mathbb{C}^{\mathcal{X}} \otimes \mathbb{C}^{\mathcal{Y}}$. For
    $k \ge 1$ the $k$-parallel query operator is $O_H^{\otimes k}$, acting
    on $k$ independent query/answer register pairs.
  \item $n$ is the number of oracle queries of the puzzle generator, $m$
    the number of oracle queries of the honest puzzle solver, and $\rho$
    the completeness error of the puzzle (the honest solver fails with
    probability at most $\rho$).
  \item $\varepsilon, \delta \in (0,1]$ are slack parameters of the
    attack: $\varepsilon$ controls its round budget and $\delta$ its
    additional loss in success probability.
  \item $r_G$ denotes the randomness of the puzzle generator, $\puz$ the
    puzzle, $\solu$ the correct solution, $\solatt$ the attacker's output.
  \item $p$ denotes a polynomial in four variables.
\end{itemize}

\section{The conjecture}\label{sec:conjectures}

\begin{definition}[rounds of quantum queries]\label{def:rounds}
Let $H \colon \mathcal{X} \to \mathcal{Y}$. A quantum oracle algorithm
$\mathsf{A}^{(\cdot)}$ \emph{asks at most $q$ queries in at most $d$ rounds}
if, on every input, its execution relative to $H$ can be written as
\[
  U_{d}\, O_H^{\otimes k_{d}}\, U_{d-1} \cdots U_{1}\,
  O_H^{\otimes k_{1}}\, U_{0}
\]
applied to a fixed initial state, followed by a measurement of a designated
output register, where the unitaries $U_0, \dots, U_{d}$ do not depend on
$H$ and the widths satisfy $k_1, \dots, k_{d} \ge 0$ and
$\sum_{i=1}^{d} k_i \le q$. Randomized algorithms are included by letting
$U_0$ prepare a uniform superposition over a randomness register that is
measured at the start. A \emph{classical-query} algorithm is the special
case in which every query register is measured in the computational basis
immediately before each application of $O_H^{\otimes k_i}$.
\end{definition}

\begin{definition}[fully quantum time-lock puzzle in the quantum random
oracle model]\label{def:qgqs-tlp}
Let $n = n(\kappa)$, $m = m(\kappa)$ and $\rho = \rho(\kappa)$. A
\emph{fully quantum time-lock puzzle with generator query complexity $n$,
solver query complexity $m$ and completeness error $\rho$} is a pair
$\Pi = (\Gen, \Sol)$ of quantum oracle algorithms such that:
\begin{itemize}
  \item $\Gen^{H}(r_G)$, on randomness $r_G$, makes at most $n$ quantum
    queries to $H$ and outputs a pair of classical strings
    $(\puz, \solu)$, the puzzle and its solution;
  \item $\Sol^{H}(\puz)$ makes at most $m$ quantum queries to $H$ and
    outputs a classical string $\solhon$;
  \item (completeness)
    \begin{multline*}
      \Pr\bigl[\solu = \solhon \;:\;
        (\puz, \solu) \leftarrow \Gen^{H}(r_G), \\
        \solhon \leftarrow \Sol^{H}(\puz)\bigr] \;\ge\; 1 - \rho,
    \end{multline*}
    over the uniform choice of $H$, the choice of $r_G$, and the
    measurements of $\Gen$ and $\Sol$.
\end{itemize}
No restriction is placed on the number of rounds in which $\Gen$ and $\Sol$
ask their queries; in particular the honest solver may need all $m$ of its
queries to be sequential, which is what a time-lock puzzle intends to
exploit.

The requirement that the puzzle $\puz$ and the solution $\solu$ be
\emph{classical} is a restriction, and the conjecture below covers only
puzzles of that form. It is the form taken by the source paper's own formal
definition (Definition 2.6, p.~11) and by its fully quantum construction
(Construction 6.13, p.~33, where $\puz = c_3$ and
$\solu = \mathsf{key}_{A_3}$), and hence the form in which its Theorem 6.2
barrier is proved. The paper explicitly contemplates, but never formalizes,
the case in which the puzzle itself is a quantum object; that case is a
strictly more general question which this statement does not address.
\end{definition}

\begin{conjecture}[breaking fully quantum time-lock puzzles with quantum
queries]\label{conj:main}
There is a fixed polynomial $p$, independent of everything below, such that
the following holds.

Let $\Pi = (\Gen, \Sol)$ be any fully quantum time-lock puzzle in the
quantum random oracle model in the sense of
Definition~\ref{def:qgqs-tlp}, with generator query complexity
$n = n(\kappa)$, solver query complexity $m = m(\kappa)$, and completeness
error $\rho = \rho(\kappa)$. Then for every
$\varepsilon, \delta \in (0,1]$ there exists a quantum oracle algorithm
$\Ext$ (allowed to depend on $\Pi$, on $\varepsilon$ and on $\delta$, and
otherwise unrestricted in its internal computation, in particular not
required to be efficient) such that:
\begin{enumerate}
  \item $\Ext$ makes at most $p(n, m, 1/\varepsilon, 1/\delta)$ quantum
    queries to $H$;
  \item those queries are arranged in at most
    $\lceil n/\varepsilon \rceil$ rounds in the sense of
    Definition~\ref{def:rounds}; and
  \item
    \begin{multline*}
      \Pr\bigl[\solatt = \solu \;:\;
        (\puz, \solu) \leftarrow \Gen^{H}(r_G), \\
        \solatt \leftarrow \Ext^{H}(\puz)\bigr]
        \;\ge\; 1 - \rho - \varepsilon - \delta,
    \end{multline*}
    where the probability is over the uniform choice of
    $H \colon \mathcal{X} \to \mathcal{Y}$, the generator randomness
    $r_G$, and the measurements of $\Ext$.
\end{enumerate}
Here $n$, $m$ and $\rho$ may be arbitrary functions of the security
parameter $\kappa$, and the statement is asserted for all sufficiently
large $\kappa$.
\end{conjecture}

\begin{corollary}[consequence, if the conjecture
holds]\label{cor:consequence}
Taking $\varepsilon = \delta = 1/4$ and $\rho \le 1/4$ would give, for
every such $\Pi$, a quantum attacker asking $\poly(n,m)$ quantum queries in
$O(n)$ rounds that recovers $\solu$ with probability at least $1/4$, so no
fully quantum time-lock puzzle in the quantum random oracle model would
have negligible soundness error against $O(n)$-round quantum attackers.
\end{corollary}

\begin{remark}[the quantitative shape is the sharp form]\label{rem:sharp}
The source paper's own words are only that ``Theorem 1.1 could potentially
be extended to the fully quantum setting using a quantum adversary'',
without pinning down the round and query bounds. The shape asserted above
--- $\lceil n/\varepsilon \rceil$ rounds,
$p(n, m, 1/\varepsilon, 1/\delta)$ queries, success
$1 - \rho - \varepsilon - \delta$ --- is therefore Theorem 3.1's shape
transplanted verbatim to the new setting rather than a bound the paper
writes down for it. An attacker achieving, say, $\poly(n)$ rounds, or
merely non-negligible rather than near-honest success, would already answer
what the authors are asking and would refute nothing here;
Conjecture~\ref{conj:main} is the sharp form.
\end{remark}

\begin{remark}[scope]\label{rem:scope}
As recorded at the end of Definition~\ref{def:qgqs-tlp}, the puzzle and its
solution are classical here. The paper observes that in the fully quantum
setting ``one can imagine the puzzle itself to be a quantum object'', and
the restriction is deliberate: it matches the paper's Definition 2.6 and is
the setting of its Theorem 6.2 barrier. It does mean that the
quantum-puzzle variant is a strictly more general question, on which
Conjecture~\ref{conj:main} is silent.
\end{remark}

\section{Bibliography}

\begin{conjurabibliography}{99}

\bibitem[AA14]{AA14}
S. Aaronson, A. Ambainis.
\emph{The need for structure in quantum speedups}.
2014. (No venue is given in the source reference list.)

\bibitem[ACC22]{ACC22}
P. Austrin, H. Chung, K.-M. Chung, S. Fu, Y.-T. Lin, M. Mahmoody.
\emph{On the impossibility of key agreements from quantum random oracles}.
Advances in Cryptology --- CRYPTO 2022: 42nd Annual International
Cryptology Conference, Proceedings, Part II, pages 165--194, Springer,
2022.

\bibitem[BDF11]{BDF11}
D. Boneh, O. Dagdelen, M. Fischlin, A. Lehmann, C. Schaffner, M. Zhandry.
\emph{Random oracles in a quantum world}.
Advances in Cryptology --- ASIACRYPT 2011, volume 7073 of Lecture Notes in
Computer Science, pages 41--69, Springer, Heidelberg, 2011.

\bibitem[CFHL21]{CFHL21}
K.-M. Chung, S. Fehr, Y.-H. Huang, T.-N. Liao.
\emph{On the compressed-oracle technique, and post-quantum security of
proofs of sequential work}.
Advances in Cryptology --- EUROCRYPT 2021, Part II, volume 12697 of Lecture
Notes in Computer Science, pages 598--629, Springer, Heidelberg, 2021.

\bibitem[MMV11]{MMV11}
M. Mahmoody, T. Moran, S. P. Vadhan.
\emph{Time-lock puzzles in the random oracle model}.
Advances in Cryptology --- CRYPTO 2011, volume 6841 of Lecture Notes in
Computer Science, pages 39--50, Springer, Heidelberg, 2011.

\bibitem[MSW20]{MSW20}
M. Mahmoody, C. Smith, D. J. Wu.
\emph{Can verifiable delay functions be based on random oracles?}
ICALP 2020: 47th International Colloquium on Automata, Languages and
Programming, volume 168 of LIPIcs, pages 83:1--83:17, Schloss Dagstuhl,
2020.

\bibitem[OSSS05]{OSSS05}
R. O'Donnell, M. Saks, O. Schramm, R. A. Servedio.
\emph{Every decision tree has an influential variable}.
46th Annual IEEE Symposium on Foundations of Computer Science (FOCS'05),
pages 31--39, IEEE, 2005.

\bibitem[RSW96]{RSW96}
R. L. Rivest, A. Shamir, D. A. Wagner.
\emph{Time-lock puzzles and timed-release crypto}.
Massachusetts Institute of Technology, Laboratory for Computer Science,
1996.

\end{conjurabibliography}

\end{document}
